  % ===========================================================
  %  main.tex  —  Sahil Saxena (11409565)
  %  A Human-in-the-Loop Framework for Alzheimer's Disease
  %  Stage Classification Using Multimodal Biomarkers
  %  University of Manchester, Department of Computer Science
  % ===========================================================

  \documentclass[12pt, a4paper]{report}

  % --- Encoding and font ---
  \usepackage[T1]{fontenc}
  \usepackage[utf8]{inputenc}
  \usepackage{microtype}

  % --- Page layout ---
  \usepackage[a4paper, top=2.5cm, bottom=2.5cm, left=3cm, right=2.5cm]{geometry}
  \usepackage{setspace}
  \onehalfspacing

  % --- Mathematics ---
  \usepackage{amsmath}
  \usepackage{amssymb}

  % --- Graphics ---
  \usepackage{graphicx}
  \usepackage{caption}
  \usepackage{subcaption}
  \usepackage{float}

  % --- Tables ---
  \usepackage{booktabs}
  \usepackage{array}
  \usepackage{tabularx}
  \usepackage{multirow}
  \usepackage{longtable}

  % --- Lists ---
  \usepackage{enumitem}

  % --- Colour ---
  \usepackage{xcolor}

  % --- Drawing (for future TikZ diagrams) ---
  \usepackage{tikz}

  % --- TOC and headings ---
  \usepackage{tocloft}
  \usepackage{titlesec}
  \usepackage{parskip}

  % --- Bibliography (natbib, author-year) ---
  \usepackage[round, authoryear, sort&compress]{natbib}

  % --- Hyperlinks (load last) ---
  \usepackage[hidelinks]{hyperref}


  % ===========================================================
  %  Heading formatting
  % ===========================================================
  \titleformat{\chapter}[block]
    {\normalfont\large\bfseries}
    {Chapter \thechapter:}{0.5em}{}[\vspace{0.5ex}\titlerule]

  \titlespacing*{\chapter}{0pt}{2ex}{2ex}

  % ===========================================================
  %  TOC formatting
  % ===========================================================
  \renewcommand{\cftchappresnum}{Chapter }
  \renewcommand{\cftchapaftersnum}{:}
  \setlength{\cftchapnumwidth}{6.5em}
  \renewcommand{\cftchapleader}{\cftdotfill{\cftdotsep}}
  \setlength{\cftbeforechapskip}{0.5em}

  % ===========================================================
  %  Convenience macros
  % ===========================================================
  \newcommand{\adni}{ADNI}
  \newcommand{\hitl}{HITL}
  \newcommand{\scd}{SCD}
  \newcommand{\mci}{MCI}
  \newcommand{\adclass}{AD}


  % ===========================================================
  \begin{document}
  % ===========================================================

  % -----------------------------------------------------------
  %  TITLE PAGE
  % -----------------------------------------------------------
  \begin{titlepage}
    \centering
    \vspace*{3cm}

    {\normalsize The University of Manchester\\
    Department of Computer Science\\[0.4em]
    Third-Year Project Report, 2025}

    \vspace{2cm}

    {\LARGE \textbf{A Human-in-the-Loop Framework for\\[0.3em]
    Alzheimer's Disease Stage Classification\\[0.3em]
    Using Multimodal Biomarkers}}

    \vspace{2cm}

    {\normalsize
    \begin{tabular}{rl}
      \textbf{Author:}     & Sahil Saxena\\[0.4em]
      \textbf{Student ID:} & 11409565\\[0.4em]
      \textbf{Supervisor:} & [Supervisor Name]
    \end{tabular}}

  \end{titlepage}


  % -----------------------------------------------------------
  %  ABSTRACT
  % -----------------------------------------------------------
  \chapter*{Abstract}
  \addcontentsline{toc}{chapter}{Abstract}

  Alzheimer's disease (AD) follows a protracted pre-clinical trajectory, passing
  through Subjective Cognitive Decline (SCD) and Mild Cognitive Impairment (MCI)
  before culminating in dementia. Accurate staging along this continuum has become
  clinically urgent with the approval of disease-modifying therapies that are most
  effective early in disease progression. Machine learning (ML) classifiers trained
  on biomarker data offer consistent, scalable staging, but their deployment in
  clinical settings is complicated by opaque predictions, variable confidence, and
  the absence of any mechanism for clinician oversight or feedback.

  This project designs and implements a Human-in-the-Loop (HITL) classification
  framework that wraps a calibrated five-model ensemble around a structured
  escalation engine, an interactive clinician dashboard, and an active-learning
  feedback loop. The ensemble combines CatBoost, Random Forest, a feedforward
  Neural Network, XGBoost, and a Support Vector Machine, trained with
  Repeated Stratified $k$-fold cross-validation ($k=5$, two repeats) on 460 patients
  from the Alzheimer's Disease Neuroimaging Initiative (ADNI). Out-of-fold
  probability outputs are aggregated via stacked generalisation, calibrated using
  sigmoid and isotonic methods, and evaluated against a strictly held-out 20\%
  test partition (116 patients). The ensemble achieves a macro one-vs-rest AUC of
  0.902 and an Expected Calibration Error of 0.042 on the out-of-fold set.

  A multi-signal escalation engine classifies each case into one of three review
  tiers based on seven uncertainty and data-quality signals, including predictive
  entropy, confidence margin, ensemble disagreement, missing-data fraction, and
  cross-modality mismatch. Conformal prediction sets provide distribution-free
  coverage guarantees alongside the point predictions. An interactive Streamlit
  dashboard exposes these signals to the clinician, alongside SHAP-based feature
  attributions and nearest-neighbour patient comparisons, and captures structured
  feedback including corrected diagnoses and confidence ratings. Clinician corrections
  feed an active-learning loop that incrementally retrains the ensemble using
  confidence-weighted sample augmentation. A simulated clinician experiment at
  85\%, 90\%, and 95\% accuracy levels demonstrates measurable cost reduction
  relative to the AI-only baseline. An ablation study, fairness audit, and
  threshold sensitivity analysis further characterise the framework's behaviour
  across demographic subgroups and operating conditions.


  % -----------------------------------------------------------
  %  TABLE OF CONTENTS
  % -----------------------------------------------------------
  \tableofcontents
  \listoffigures
  \listoftables


  % ===========================================================
  \chapter{Introduction}
  \label{chap:introduction}
  % ===========================================================

  % -----------------------------------------------------------
  \section{Motivation and Clinical Context}
  \label{sec:motivation}
  % -----------------------------------------------------------

  Alzheimer's disease (AD) is the most prevalent form of dementia, accounting for
  an estimated 60--70\% of all cases worldwide \citep{who2023dementia}. The disease
  follows a progressive neurodegenerative trajectory that begins years before clinical
  symptoms become apparent, passing through identifiable preclinical and prodromal
  stages before culminating in irreversible cognitive decline \citep{jack2018nia}.
  Once a patient receives a diagnosis of AD, the scope for therapeutic intervention
  is severely limited; current pharmacological treatments offer only modest symptomatic
  relief and cannot reverse the underlying neuronal damage. The consequences extend
  well beyond the individual patient: family members assume substantial caregiving
  responsibilities, healthcare systems face escalating costs, and society at large
  contends with a condition for which no cure yet exists \citep{wimo2017costs}.

  What makes the disease particularly challenging from a clinical standpoint is the
  continuum along which it develops. Before a patient meets the criteria for an AD
  diagnosis, they typically pass through two earlier stages. Subjective Cognitive
  Decline (SCD) describes the phase in which individuals report perceived worsening
  of their cognitive abilities yet perform within normal ranges on standardised
  objective tests \citep{jessen2014scd}. A proportion of these individuals
  subsequently progress to Mild Cognitive Impairment (MCI), a stage characterised by
  measurable deficits in one or more cognitive domains that nonetheless fall short of
  the threshold for dementia \citep{albert2011mci}. Not all SCD patients convert to
  MCI, and not all MCI patients convert to AD---some remain stable, while others
  develop non-Alzheimer's dementias. For those on the amnestic trajectory, however,
  early identification opens a critical window for intervention.

  This window has acquired renewed urgency with the emergence of disease-modifying
  therapies (DMTs). Recent trials of anti-amyloid monoclonal antibodies---lecanemab
  and donanemab---have demonstrated statistically significant reductions in the rate
  of cognitive decline among patients in the early symptomatic stages of AD
  \citep{vandyck2023lecanemab, sims2023donanemab}. Both agents are most effective
  when administered before extensive neurodegeneration has occurred, making accurate
  staging along the SCD--MCI--AD continuum a direct clinical prerequisite for
  treatment eligibility.

  Yet accurate staging remains difficult. The boundaries between SCD, MCI, and AD
  are inherently fuzzy, clinical assessments are subject to inter-rater variability,
  and the biomarker profiles that distinguish one stage from another frequently
  overlap \citep{edmonds2015mci}. Machine learning (ML) methods offer a promising
  route to more consistent, data-driven classification by learning complex patterns
  across multiple biomarker modalities simultaneously. However, deploying ML models
  in clinical settings raises a distinct set of concerns: clinicians are unlikely to
  accept a fully automated diagnostic system for a condition as consequential as
  Alzheimer's disease, particularly when model confidence is low or the patient
  profile is atypical. This tension between the pattern-recognition strengths of
  automated classifiers and the contextual judgement of experienced clinicians
  motivates a Human-in-the-Loop (HITL) approach, in which the AI system handles
  routine cases autonomously while escalating uncertain or high-risk cases for
  expert review.

  % -----------------------------------------------------------
  \section{Problem Statement}
  \label{sec:problem}
  % -----------------------------------------------------------

  Existing ML approaches to Alzheimer's stage classification typically operate in
  isolation: a model receives patient data, produces a prediction, and that
  prediction is either accepted or discarded with no mechanism for structured clinical
  oversight or feedback. This produces two compounding failure modes. First,
  misclassifications go uncorrected, meaning patients may be assigned to the wrong
  stage with downstream consequences for treatment eligibility and clinical
  management. Second, the model never benefits from clinician expertise; errors that
  a domain expert could readily identify are repeated indefinitely because no feedback
  pathway exists.

  This project addresses both failure modes by designing and implementing a HITL
  framework that wraps an ensemble classifier for three-class staging (SCD, MCI, AD)
  using biomarker data from the Alzheimer's Disease Neuroimaging Initiative (ADNI).
  The framework introduces an escalation engine that identifies cases where the
  ensemble is most likely to be wrong, routes those cases to a clinician for review,
  captures structured feedback, and uses that feedback to iteratively improve the
  underlying model through active learning. The central question is whether this
  selective human--AI collaboration achieves materially better classification accuracy
  and lower clinical risk than either the AI system or a simulated clinician
  operating alone.

  % -----------------------------------------------------------
  \section{Aims and Objectives}
  \label{sec:aims}
  % -----------------------------------------------------------

  \subsection*{Aim}

  To develop a HITL classification framework that combines a calibrated ensemble ML
  system with structured clinician oversight to improve the accuracy, safety, and
  interpretability of Alzheimer's disease stage classification across SCD, MCI, and AD.

  \subsection*{Objectives}

  \begin{enumerate}[leftmargin=*, label=\textbf{O\arabic*.}]

    \item \textbf{Calibrated Ensemble Classification Pipeline}
    \begin{enumerate}[label=(\alph*)]
      \item Build a five-model calibrated ensemble (CatBoost, Random Forest, Neural
      Network, XGBoost, SVM) trained on ADNI biomarker data, using out-of-fold
      predictions from Repeated Stratified $k$-fold cross-validation to avoid data
      leakage.
      \item \textit{Success criterion:} The ensemble should outperform each individual
      base learner on macro F1 and AUC, with performance reported as mean and standard
      deviation across folds.
    \end{enumerate}

    \item \textbf{Multi-Signal Escalation Engine}
    \begin{enumerate}[label=(\alph*)]
      \item Design a set of uncertainty and data-quality signals (predictive
      confidence, decision margin, entropy, ensemble disagreement, missing-data
      fraction, cross-modality mismatch, SHAP instability) and use them to partition
      each case into one of three review tiers: AI-Autonomous, AI-Assisted, and
      Clinician-Mandatory.
      \item \textit{Success criterion:} The escalation policy should correctly capture
      the majority of misclassified cases while maintaining an interpretable and
      adjustable escalation rate.
    \end{enumerate}

    \item \textbf{Clinician Dashboard}
    \begin{enumerate}[label=(\alph*)]
      \item Develop an interactive Streamlit-based interface that allows clinicians to
      review escalated cases, inspect AI reasoning through SHAP feature attributions,
      examine patient biomarker values relative to reference ranges, and submit
      structured feedback including a corrected diagnosis and a confidence rating on a
      1--5 scale.
      \item \textit{Success criterion:} All information necessary for an informed
      clinical decision is accessible on a single screen, with feedback captured in a
      structured, auditable SQLite log.
    \end{enumerate}

    \item \textbf{Active-Learning Feedback Loop}
    \begin{enumerate}[label=(\alph*)]
      \item Implement a mechanism whereby clinician corrections are accumulated into a
      feedback pool and used to periodically retrain the ensemble, with clinician
      confidence ratings converted to sample weights.
      \item \textit{Success criterion:} Classification accuracy on a held-out test set
      should improve over successive feedback rounds relative to the AI-only baseline.
    \end{enumerate}

    \item \textbf{Experimental Evaluation}
    \begin{enumerate}[label=(\alph*)]
      \item Conduct a controlled comparison of AI-only and HITL classification under
      simulated clinician conditions at accuracy levels of 85\%, 90\%, and 95\%,
      supported by bootstrap confidence intervals, an escalation ablation study,
      threshold sensitivity analysis, cost--benefit analysis under a clinically
      motivated error-cost matrix, and a fairness audit across age and gender
      subgroups.
      \item \textit{Success criterion:} The HITL system should demonstrate lower
      cost-weighted error than the AI-only baseline, and the fairness audit should
      quantify any demographic disparities in escalation rates.
    \end{enumerate}

  \end{enumerate}

  % -----------------------------------------------------------
  \section{Evaluation Strategy}
  \label{sec:evaluation-strategy}
  % -----------------------------------------------------------

  Evaluation operates at three levels. At the \emph{model level}, each base learner
  and the calibrated ensemble are assessed using accuracy, macro-averaged F1, per-class
  precision and recall, one-vs-rest AUC, Expected Calibration Error (ECE), and
  bootstrap 95\% confidence intervals, all derived from out-of-fold predictions to
  prevent leakage. At the \emph{escalation level}, an ablation study removes each
  uncertainty signal in turn to quantify its marginal contribution, and a threshold
  sensitivity analysis characterises the accuracy--review-rate trade-off. At the
  \emph{system level}, the HITL pipeline is compared against the AI-only baseline
  under a clinically motivated cost matrix that penalises dangerous misclassifications
  (e.g.\ classifying an AD patient as SCD) more heavily than near-stage errors.
  Conformal prediction provides set-valued predictions with finite-sample
  coverage guarantees \citep{angelopoulos2023conformal}. A fairness audit examines
  whether escalation rates and accuracy differ across subgroups defined by age and
  gender, using demographic parity and equalised odds as metrics
  \citep{mehrabi2021survey}.

  % -----------------------------------------------------------
  \section{Report Structure}
  \label{sec:structure}
  % -----------------------------------------------------------

  Chapter~\ref{chap:clinical-background} provides the clinical foundation, covering
  Alzheimer's disease pathology, the SCD--MCI--AD staging continuum, the challenges
  of clinical diagnosis, and the structure of the ADNI dataset used in this project.
  Chapter~\ref{chap:litreview} reviews the technical literature, spanning machine
  learning for AD classification, ensemble methods, Human-in-the-Loop systems,
  uncertainty quantification and conformal prediction, explainable AI, and active
  learning. It concludes by positioning this project within the research landscape.
  Chapter~\ref{chap:design} describes the overall system design and high-level
  architecture, motivating the key design decisions. Subsequent chapters present the
  detailed implementation, experimental results, critical evaluation, and conclusions
  with directions for future work.

  The source code, pre-trained model artefacts, and generated evaluation reports are
  maintained in a version-controlled repository. All experimental results reported
  in this dissertation are reproducible by executing the provided pipeline scripts
  from a clean environment.


  % ===========================================================
  \chapter{Clinical Background}
  \label{chap:clinical-background}
  % ===========================================================

  % -----------------------------------------------------------
  \section{Alzheimer's Disease Pathology}
  \label{sec:ad-pathology}
  % -----------------------------------------------------------

  Alzheimer's disease is a progressive neurodegenerative disorder characterised at
  the molecular level by two principal neuropathological hallmarks: extracellular
  deposits of amyloid-$\beta$ (A$\beta$) peptide in the form of neuritic plaques, and
  intraneuronal accumulations of hyperphosphorylated tau protein in the form of
  neurofibrillary tangles (NFTs). The \emph{amyloid cascade hypothesis} proposes that
  aberrant cleavage of amyloid precursor protein (APP) by $\beta$- and
  $\gamma$-secretases generates A$\beta_{42}$ fragments that misfold and aggregate,
  triggering a cascade of downstream events including neuroinflammation, oxidative
  stress, synaptic dysfunction, and ultimately the phosphorylation and aggregation of
  tau \citep{mckhann2011diagnosis}.

  The temporal ordering of these pathological events is well-characterised and forms
  the basis of the \emph{A/T/(N) biomarker classification system} proposed by
  \citet{jack2018nia} for the National Institute on Aging and Alzheimer's Association
  (NIA-AA). Under this framework, biomarkers are grouped into three categories:
  \textbf{A} (amyloid pathology), measured by cerebrospinal fluid (CSF)
  A$\beta_{42}$, the A$\beta_{42}$/A$\beta_{40}$ ratio, or amyloid PET; \textbf{T}
  (tau pathology), measured by CSF phosphorylated tau (p-tau) or tau PET; and
  \textbf{(N)} (neurodegeneration), measured by CSF total tau (t-tau), FDG-PET
  hypometabolism, or structural MRI atrophy. Critically, amyloid pathology is
  detectable approximately 15--20 years before the onset of dementia symptoms,
  providing a substantial window for biomarker-guided intervention.

  Recent advances in plasma-based biomarkers---particularly plasma
  neurofilament light chain (NfL) and plasma p-tau217---have made non-invasive
  biomarker screening increasingly feasible \citep{zetterberg2021plasma}. NfL is a
  marker of axonal injury that rises across a range of neurodegenerative conditions,
  while elevated plasma p-tau reliably distinguishes AD pathology from other
  dementia subtypes. The availability of these plasma markers in large-scale cohort
  studies such as ADNI has opened new possibilities for computationally tractable,
  multi-modal classification without requiring lumbar puncture.

  % -----------------------------------------------------------
  \section{The AD Continuum: SCD, MCI, and Dementia}
  \label{sec:ad-continuum}
  % -----------------------------------------------------------

  The contemporary understanding of Alzheimer's disease frames it not as an acute
  condition but as a decades-long biological continuum that unfolds in clinically
  distinct stages. The first recognisable stage, \emph{Subjective Cognitive Decline}
  (SCD), is defined by self-reported perception of worsening memory or cognitive
  function in the absence of detectable objective impairment on standardised
  neuropsychological tests \citep{jessen2014scd}. Although SCD individuals are
  cognitively within normal limits, a substantial subset carry early AD pathology:
  studies have shown that elevated CSF tau-to-amyloid ratios, amyloid-PET positivity,
  and accelerated hippocampal volume loss are already detectable at this stage.
  Identifying which SCD individuals are on a trajectory to conversion is the primary
  clinical challenge the DMT era now demands.

  The second stage, \emph{Mild Cognitive Impairment} (MCI), is operationally defined
  as objective impairment on cognitive testing---typically one to two standard
  deviations below age-adjusted norms---combined with functional independence in daily
  activities \citep{albert2011mci, petersen2004mci}. MCI due to Alzheimer's
  disease (\emph{amnestic MCI}) is the most clinically significant subtype, carrying
  an annual conversion rate to dementia of approximately 10--15\% compared with
  1--2\% in cognitively normal individuals. Non-amnestic MCI, in contrast, is more
  likely to progress to other forms of dementia and may not benefit from
  anti-amyloid therapies.

  \emph{Alzheimer's dementia} is diagnosed when cognitive deficits are sufficiently
  severe to interfere with daily functioning \citep{mckhann2011diagnosis}. By this
  stage, widespread synaptic loss, cortical atrophy (particularly in the entorhinal
  cortex, hippocampus, and temporal lobes), and significant NFT burden are
  established. The clinical utility of AD diagnostics therefore pivots increasingly
  towards the pre-dementia stages, where intervention is most likely to be effective.

  For DMT eligibility, this staging precision is critical. Lecanemab was trialled in
  patients with early AD (MCI or mild dementia with amyloid confirmation), achieving
  a 27\% reduction in clinical decline over 18 months \citep{vandyck2023lecanemab}.
  Donanemab similarly targeted early symptomatic patients, achieving a 35\%
  slowing of disease progression in the low/medium tau subgroup \citep{sims2023donanemab}.
  Both results underscore that the therapeutic window is narrow and that accurate
  pre-selection---distinguishing SCD converters, amnestic MCI, and established
  AD from non-converters and non-AD conditions---translates directly to treatment
  benefit.

  % -----------------------------------------------------------
  \section{Clinical Diagnosis and the Role of AI}
  \label{sec:clinical-diagnosis}
  % -----------------------------------------------------------

  Current clinical diagnosis of Alzheimer's disease stages relies on the convergence
  of evidence from multiple modalities: standardised cognitive assessments such as
  the Mini-Mental State Examination (MMSE), Montreal Cognitive Assessment (MoCA),
  and the Clinical Dementia Rating (CDR); structural neuroimaging to assess cortical
  atrophy and hippocampal volume loss; CSF biomarkers for A$\beta$ and tau; and
  increasingly, plasma biomarkers for NfL and p-tau217. A definitive ante-mortem
  diagnosis still requires biomarker confirmation in most research-grade protocols.

  This multi-modal assessment process is resource-intensive, time-consuming, and
  requires specialised expertise. Moreover, the classification boundaries between
  SCD, MCI, and AD are not hard thresholds but probabilistic gradients, meaning that
  clinicians must exercise considerable judgement when interpreting borderline profiles.
  Studies have shown that conventional MCI criteria are susceptible to false-positive
  diagnostic errors, and that agreement rates between specialists can be moderate
  \citep{edmonds2015mci}. This variability introduces both under-diagnosis---missing
  patients who would benefit from early treatment---and over-diagnosis, exposing
  patients to unnecessary follow-up or, in the case of DMTs, treatments with
  significant side-effect profiles.

  Supervised machine learning offers a complementary path. Models trained on large,
  curated biomarker datasets can identify combinations of features that predict staging
  labels more reliably than individual biomarkers in isolation, and can generalise
  these patterns consistently without fatigue or inter-rater drift
  \citep{rathore2017review}. However, applying ML to clinical decision support
  introduces well-documented challenges: black-box predictions erode clinician trust,
  high-confidence errors can propagate uncorrected, and there is no inherent mechanism
  for incorporating domain expertise or patient context. A HITL framework addresses
  these shortcomings by designating the ML system as a first-line screener that
  identifies and quantifies its own uncertainty, and reserving final authority for
  clinicians on the cases where that uncertainty is greatest \citep{topol2019highperformance}.

  % -----------------------------------------------------------
  \section{The ADNI Dataset}
  \label{sec:adni}
  % -----------------------------------------------------------

  The Alzheimer's Disease Neuroimaging Initiative (ADNI) is a longitudinal,
  multi-site observational study established in 2004, currently in its fourth phase
  (ADNI4), with the primary goal of validating biomarkers for use as clinical trial
  endpoints in AD studies \citep{weiner2017adni}. ADNI recruits participants across
  three diagnostic categories---cognitively normal, MCI, and AD---and collects
  repeated measurements of MRI volumetrics, amyloid and tau PET imaging, CSF
  biomarkers, plasma biomarkers, genetic markers (including APOE genotyping), and a
  battery of neuropsychological assessments.

  The dataset used in this project is derived from the ADNI biomarker archive and
  contains \textbf{576 patients} mapped to three diagnostic labels: SCD, MCI, and AD.
  Following the convention of the preceding student codebase from which this project
  extends, the label SCD subsumes the cognitively normal and subjective complaint
  categories. The feature set, after merging longitudinal records and applying the
  preprocessing pipeline described in Chapter~\ref{chap:design}, spans four
  modalities:

  \begin{itemize}[nosep]
    \item \textbf{CSF biomarkers:} total tau (TAU), phosphorylated tau (PTAU),
      amyloid-$\beta_{1\text{--}42}$ (ABETA), and the A$\beta_{42}$/A$\beta_{40}$
      ratio (AB4240). CSF biomarkers are available for approximately 65--66\% of
      patients, reflecting the invasiveness of lumbar puncture.
    \item \textbf{Plasma biomarkers:} neurofilament light chain (PLASMA\_NFL) and
      plasma p-tau (PLASMATAU), available for 100\% and 99.5\% of patients
      respectively, consistent with the lower barrier of blood-based testing.
    \item \textbf{Cognitive assessment:} MMSE score, available for all 576 patients.
    \item \textbf{Demographic features:} age (AGE), sex (PTGENDER), and years of
      education (PTEDUCAT), all fully observed.
  \end{itemize}

  Five engineered ratio features (e.g.\ TAU/PTAU, ABETA/TAU) augment the raw
  biomarkers, yielding 16 candidate features in total. The high missingness rate on
  CSF variables is clinically expected and is handled within the pipeline using
  distance-weighted $k$-NN imputation, as described in Chapter~\ref{chap:design}.
  Crucially, imputation is performed within each cross-validation fold to prevent
  information leakage from held-out validation samples into imputed training values.


  % ===========================================================
  \chapter{Technical Background and Literature Review}
  \label{chap:litreview}
  % ===========================================================

  % -----------------------------------------------------------
  \section{Machine Learning for Alzheimer's Disease Classification}
  \label{sec:ml-ad}
  % -----------------------------------------------------------

  The application of supervised machine learning to Alzheimer's disease classification
  has a two-decade history, evolving from early support vector machine (SVM) classifiers
  applied to structural MRI to contemporary multi-modal deep learning pipelines.
  \citet{kloppel2008automatic} demonstrated that a linear SVM operating on voxel-based
  morphometry (VBM) features could discriminate AD from cognitively normal individuals
  with accuracy comparable to expert radiologists, establishing the proof of concept
  for automated biomarker-based classification. The domain subsequently diversified:
  random forests, gradient boosting machines, and multi-layer perceptrons have all
  been applied to tabular biomarker data, while convolutional neural networks and
  graph neural networks have been applied to imaging modalities.

  \citet{rathore2017review} provide a systematic review of neuroimaging-based
  classification studies for AD and its prodromal stages, documenting the consistent
  advantage of multi-modal feature sets over single-modality inputs. The review
  identifies class imbalance, small sample sizes, and heterogeneous acquisition
  protocols across sites as the three primary methodological challenges in the field.
  For tabular biomarker datasets such as ADNI, the additional challenge of variable
  missingness---arising from the differing availability of CSF, plasma, and imaging
  data across participants---necessitates careful imputation strategies that are
  integrated into the evaluation pipeline rather than applied globally.

  Deep learning approaches have achieved strong results when sufficiently large imaging
  datasets are available. \citet{ding2019deep} trained a convolutional network on
  $^{18}$F-FDG PET scans and demonstrated early detection of AD pathology up to six
  years before clinical diagnosis, highlighting the longitudinal potential of imaging
  biomarkers. However, for smaller tabular cohorts such as the one used in this
  project (576 patients, 12 selected features), deep learning offers no systematic
  advantage over well-regularised ensemble methods, and interpretability considerations
  favour approaches with tractable feature attribution.

  Multi-class staging (SCD vs.\ MCI vs.\ AD) is intrinsically harder than binary
  detection due to the fuzzy boundaries between adjacent classes, particularly between
  SCD and MCI. Models evaluated on three-class tasks typically exhibit lower MCI
  performance than SCD or AD performance, consistent with MCI representing the most
  heterogeneous and diagnostically uncertain category. This pattern motivates
  uncertainty-aware classification: a model that is well-calibrated and capable of
  expressing its own ambiguity about MCI-adjacent cases is more clinically useful than
  one that issues overconfident predictions.

  % -----------------------------------------------------------
  \section{Ensemble Methods and Model Stacking}
  \label{sec:ensemble}
  % -----------------------------------------------------------

  Ensemble learning improves predictive performance by combining the outputs of
  multiple base learners, reducing variance through diversity \citep{breiman2001random}.
  The three principal paradigms---bagging, boosting, and stacking---are all relevant
  to this project.

  \emph{Bagging} (Bootstrap Aggregating) trains independent models on bootstrap
  samples of the training data and averages their predictions. Random Forest
  \citep{breiman2001random} combines bagging with feature subsampling at each split,
  producing a robust ensemble of decision trees that generalises well on tabular data
  and is naturally resistant to overfitting. \emph{Boosting} constructs an additive
  model sequentially, where each learner corrects the residual errors of its
  predecessors. XGBoost \citep{chen2016xgboost} achieves state-of-the-art results on
  a wide range of structured prediction tasks through second-order gradient
  approximations and L1/L2 regularisation. CatBoost \citep{prokhorenkova2018catboost}
  extends this framework with ordered boosting, which mitigates target leakage in
  the gradient computation. Both gradient boosting variants are included in the
  ensemble specifically because their inductive biases complement those of the
  variance-reducing Random Forest and the non-linear Neural Network.

  \emph{Stacked generalisation} \citep{wolpert1992stacked} trains a meta-learner
  (Level-1 model) on the out-of-fold probability outputs of the base models rather
  than their raw features, allowing the meta-learner to learn optimal linear
  combinations of the base models' probability estimates. Using out-of-fold predictions
  to construct the Level-1 training set is essential: training the meta-learner on the
  same data as the base models would cause it to upweight models that overfit,
  producing overly optimistic estimates. In this project, a logistic regression
  meta-learner is trained on a 15-dimensional feature vector (five models times three
  class probabilities) and its performance is compared against simple probability
  averaging; the better-performing combination is selected for downstream use.

  \emph{Probability calibration} is a prerequisite for a credible escalation engine:
  if a model that outputs a confidence of 0.9 is actually correct only 75\% of the
  time, the confidence threshold used to trigger human review is meaningless.
  \citet{platt1999probabilistic} proposed fitting a sigmoid function (logistic
  regression) to post-hoc map raw model scores to calibrated probabilities;
  \citet{niculescu2005predicting} extended this comparison to include isotonic
  regression calibration. Both methods are implemented here using a cross-validated
  fitting procedure that avoids the bias introduced by calibrating and evaluating on
  the same data. Expected Calibration Error (ECE) is used as the selection criterion
  between calibration methods.

  % -----------------------------------------------------------
  \section{Human-in-the-Loop Machine Learning}
  \label{sec:hitl}
  % -----------------------------------------------------------

  A Human-in-the-Loop (HITL) system is one in which human judgement is integrated
  into the machine learning pipeline in a principled and structured way, rather than
  being treated as an external intervention after the fact
  \citep{monarch2021hitl, amershi2014power}. The defining characteristic of a HITL
  system is that human feedback is \emph{anticipated by design}: the system identifies
  where human input is most valuable, presents the relevant information efficiently,
  and closes the loop by using the feedback to update the model.

  The distinction between HITL and human-in-the-loop decision support is important.
  In the latter, a human reviews every AI output and has the option to override it; in
  HITL, the system selectively routes only a subset of cases to the human, using an
  escalation policy based on the system's own uncertainty or the clinical stakes of
  the decision. This selective routing is economically necessary in any clinical
  deployment at scale: reviewing every prediction defeats the efficiency advantage
  of automation and introduces human fatigue.

  In the medical AI literature, the case for HITL deployment is compelling.
  \citet{topol2019highperformance} argues that AI systems perform best not when they
  replace clinicians but when they augment clinical capacity by handling routine,
  high-confidence cases and directing clinician attention toward complex, uncertain ones.
  \citet{rajpurkar2022ai} survey the landscape of AI in medicine and conclude that
  most high-performing medical AI systems still require some form of human oversight
  to be safe and clinically acceptable. For AD staging specifically, the combination
  of high diagnostic stakes, imperfect biomarker sensitivity, and the direct link to
  DMT eligibility makes unmoderated fully-automated classification inappropriate with
  current technology.

  In this project, the escalation policy operates at three tiers.
  \textbf{AI-Autonomous} cases are those where all uncertainty signals fall below
  their respective thresholds and no data-quality flags are raised; the system's
  prediction is accepted without review. \textbf{AI-Assisted} cases are those where
  one or more uncertainty signals are elevated, prompting a clinician review that is
  informed by the AI recommendation and SHAP explanations. \textbf{Clinician-Mandatory}
  cases are those where critical data is missing or cross-modality biomarker
  inconsistencies are detected, and the case is referred to the clinician regardless
  of prediction confidence. This tiered structure reflects the clinical principle of
  graded risk management: not all uncertain cases require the same level of human
  involvement.

  % -----------------------------------------------------------
  \section{Uncertainty Quantification and Conformal Prediction}
  \label{sec:uq}
  % -----------------------------------------------------------

  Meaningful uncertainty quantification is the foundation on which the escalation
  engine rests. For a probabilistic classifier, the most commonly used scalar
  uncertainty measures are \emph{maximum predicted probability} (confidence),
  \emph{decision margin} (the gap between the top-1 and top-2 class probabilities),
  and \emph{predictive entropy}:
  \begin{equation}
    H(\mathbf{p}) = -\sum_{k=1}^{K} p_k \log p_k,
    \label{eq:entropy}
  \end{equation}
  where $p_k$ is the predicted probability for class $k$ and $K=3$ in this project.
  Entropy reaches its maximum of $\log K$ when the model distributes probability
  uniformly across classes, and reaches zero for a perfectly confident prediction.
  These scalar measures are complementary: a low-confidence prediction may still have
  a large margin (if the probability mass is concentrated on just two classes), while
  a high-entropy prediction captures three-way ambiguity.

  These point-wise uncertainty measures are useful but lack formal statistical
  guarantees. \emph{Conformal prediction} \citep{shafer2008tutorial,
  angelopoulos2023conformal} provides a distribution-free framework for constructing
  \emph{prediction sets}---sets of candidate labels that are guaranteed to contain the
  true label with pre-specified probability $1-\alpha$, regardless of the underlying
  data distribution and the quality of the model's probability estimates. A conformal
  predictor computes a \emph{nonconformity score} for each calibration sample (in this
  project, the complement of the predicted probability for the true class), estimates
  the $(1-\alpha)$-quantile $\hat{q}$ of the empirical distribution of these scores,
  and includes class $k$ in the prediction set for a test sample if its predicted
  probability exceeds $1 - \hat{q}$.

  The guarantee is finite-sample and assumption-free: if calibration and test data are
  exchangeable, the marginal coverage is at least $1 - \alpha$ over random draws of
  the calibration set. For clinical decision support, this is a meaningful property:
  rather than claiming ``the model predicts MCI with 68\% confidence'', the system can
  state ``the model is $\geq\!95\%$ certain the correct diagnosis is in the set
  $\{\text{SCD}, \text{MCI}\}$''. Prediction sets that contain all three classes are
  flagged as maximally uncertain and are strong candidates for escalation.

  % -----------------------------------------------------------
  \section{Explainable AI}
  \label{sec:xai}
  % -----------------------------------------------------------

  Clinician trust in AI-assisted diagnosis is contingent not only on prediction
  accuracy but on the intelligibility of the reasoning behind each prediction.
  Explainable AI (XAI) methods attempt to bridge this gap by attributing contributions
  to individual input features for a given prediction.

  SHAP (SHapley Additive exPlanations) \citep{lundberg2017shap} provides a
  principled, game-theoretic framework for feature attribution based on Shapley values
  from cooperative game theory. For a prediction $f(\mathbf{x})$, each feature $i$
  receives an attribution $\phi_i$ such that:
  \begin{equation}
    f(\mathbf{x}) = \phi_0 + \sum_{i=1}^{p} \phi_i,
    \label{eq:shap}
  \end{equation}
  where $\phi_0 = \mathbb{E}[f(\mathbf{X})]$ is the expected model output and
  $\sum_i \phi_i$ is the total deviation from this baseline. The Shapley value for
  feature $i$ is computed by averaging marginal contributions over all possible
  feature subsets, guaranteeing efficiency (attributions sum to the prediction),
  symmetry (identically contributing features receive equal attribution), and
  missingness (absent features contribute zero).

  For tree-based models such as CatBoost and Random Forest, SHAP's TreeExplainer
  computes exact Shapley values in polynomial time \citep{lundberg2017shap}. This
  efficiency is critical for a real-time clinical dashboard where explanations must be
  generated on-the-fly for each reviewed patient. In the HITL dashboard, SHAP
  waterfall plots show the top contributing features for the patient's predicted class,
  enabling clinicians to identify whether the model's reasoning aligns with clinical
  expectations---for example, whether an AD prediction is driven by elevated tau and
  low amyloid-$\beta$, or by a less clinically intuitive feature combination.

  The \emph{SHAP instability} signal used in the escalation engine computes the
  variance of SHAP attribution vectors across the base models, capturing cases where
  different models attribute importance to different features for the same patient.
  High instability suggests that the ensemble's prediction is not grounded in a
  consistent biomarker pattern across learners, which is a meaningful warning even
  when the ensemble's predicted class is unanimous.

  % -----------------------------------------------------------
  \section{Active Learning and Feedback Loops}
  \label{sec:active-learning}
  % -----------------------------------------------------------

  Active learning is a semi-supervised learning paradigm in which the model identifies
  which unlabelled examples would most improve its performance if labelled, and
  selectively queries an oracle (typically a human annotator) for those labels
  \citep{settles2009active}. In the context of clinical AI, the oracle is the
  clinician, and the ``label queries'' are the cases escalated for review. The key
  insight is that not all corrections are equally informative: a clinician who
  overrides a confident, correct prediction adds no signal, while a correction on a
  genuinely uncertain boundary case has high epistemic value.

  The most commonly used active learning query strategy is \emph{uncertainty
  sampling}, which prioritises examples with the highest predictive entropy or lowest
  confidence. An alternative, \emph{query by committee} \citep{settles2009active},
  submits examples on which a committee of models disagrees most strongly, which
  aligns naturally with the ensemble disagreement signal already computed for
  escalation purposes.

  In this project, the feedback loop implements a confidence-weighted retraining
  strategy that extends standard uncertainty sampling. Clinician corrections are
  accumulated into a feedback pool ranked by the escalation risk score (a composite
  of the uncertainty signals). After each batch of corrections, the ensemble is
  retrained on the union of the original training set and the feedback pool, with
  each corrected sample weighted by a function of the clinician's self-reported
  confidence on a 1--5 scale. This weighting ensures that high-confidence corrections
  exert greater influence on the retrained model than uncertain ones, providing a
  principled mechanism for handling annotation noise.

  % -----------------------------------------------------------
  \section{Summary and Research Positioning}
  \label{sec:positioning}
  % -----------------------------------------------------------

  The preceding sections have characterised a research landscape in which several
  active research threads---biomarker-based AD classification, ensemble learning with
  calibration, HITL systems, uncertainty quantification, and active learning---have
  largely developed in isolation. Existing ML approaches to AD staging predominantly
  optimise predictive accuracy on held-out test sets without considering the
  downstream clinical workflow in which predictions will be consumed. HITL frameworks
  in the medical AI literature have been studied primarily in the context of image
  annotation and radiology, where the modality and the annotation process are
  qualitatively different from tabular biomarker data.

  This project occupies a specific gap: it builds a complete, end-to-end HITL
  pipeline for multi-class tabular biomarker data in a domain where accurate
  pre-treatment staging is clinically urgent. The contributions are as follows.
  First, it demonstrates stacked generalisation with probability calibration for
  three-class AD staging on ADNI data, evaluated via leak-free out-of-fold
  cross-validation. Second, it implements a multi-signal escalation engine that
  combines seven distinct uncertainty and data-quality indicators under a principled
  three-tier policy, and formally characterises the contribution of each signal
  through ablation. Third, it provides conformal prediction sets as a rigorous,
  assumption-free complement to point-wise uncertainty estimates, offering finite-sample
  coverage guarantees that are directly communicable to clinicians. Fourth, it
  integrates all components into an interactive Streamlit dashboard with real-time
  SHAP explanations, a structured feedback capture mechanism, and an active-learning
  retraining loop. Fifth, it quantifies the clinical benefit of human oversight
  through a cost--benefit analysis under a clinically motivated error-cost matrix and
  a fairness audit across demographic subgroups. Together, these components constitute
  a more complete HITL framework for AD staging than has been reported in the
  existing literature.


  % ===========================================================
  \chapter{System Design and Architecture}
  \label{chap:design}
  % ===========================================================

  % -----------------------------------------------------------
  \section{High-Level Architecture}
  \label{sec:architecture}
  % -----------------------------------------------------------

  The system is organised as a five-stage sequential pipeline, each stage consuming
  the outputs of the previous one and producing artefacts consumed by the next. This
  linear structure was chosen over a more tightly coupled architecture because it
  allows each component to be developed, tested, and replaced independently, and
  because the clinical workflow it mirrors is itself sequential: data arrives, the
  model classifies, uncertain cases are escalated, a clinician reviews, and feedback
  returns to the model.

  \begin{figure}[htbp]
    \centering
    % === TODO: insert TikZ or high-resolution architecture diagram here ===
    \fbox{\parbox{0.88\textwidth}{%
      \centering\vspace{4cm}
      \textit{[Architecture diagram --- to be inserted]}\\[0.3em]
      \textit{Five stages: Data Preprocessing $\rightarrow$ Ensemble Classification
      $\rightarrow$ Escalation Engine $\rightarrow$ Clinician Dashboard
      $\rightarrow$ Feedback Loop}
      \vspace{4cm}}}
    \caption{High-level architecture of the HITL framework. Data flow proceeds
    left to right through five principal components. Dashed return arrows indicate
    the active-learning path by which clinician corrections re-enter the training set.}
    \label{fig:architecture}
  \end{figure}

  \paragraph{Stage 1 — Data Preprocessing.}
  Raw ADNI records from the merged CSV are ingested by the preprocessing pipeline
  (\texttt{Preprocess.py}). Non-predictive identifiers (\texttt{RID},
  \texttt{VISCODE}) are removed. Censored biomarker strings (e.g.\ \texttt{>1300}
  for TAU, \texttt{<200} for ABETA) are replaced with their numeric boundaries before
  conversion to float. Five ratio features are engineered from raw biomarker pairs
  (e.g.\ TAU/PTAU, ABETA/TAU) to capture the relative relationships that clinical
  literature identifies as diagnostically informative. Outliers in continuous features
  are capped at $Q_1 - 1.5\,\text{IQR}$ and $Q_3 + 1.5\,\text{IQR}$ before any
  missing-value treatment. The diagnostic label is encoded as an integer class
  $\{0\,{=}\,\text{SCD},\; 1\,{=}\,\text{MCI},\; 2\,{=}\,\text{AD}\}$.

  Within the model training script (\texttt{ModelsFinal.py}), a stratified 80/20
  split isolates a strictly held-out test partition (116 patients) before any
  cross-validation begins. Within each cross-validation fold, a $k$-NN imputer
  ($k=9$, distance-weighted) is fitted on the fold's training split and applied to
  the validation split, preventing target leakage through the imputation step.
  SelectKBest with mutual information scoring then selects the 12 most informative
  features, and a RobustScaler normalises them. Both the selector and the scaler are
  fitted within the fold and applied only to the corresponding validation data.

  \paragraph{Stage 2 — Ensemble Classification.}
  Five base learners are trained within each cross-validation fold: CatBoost
  (gradient-boosted trees with ordered boosting), Random Forest (bagged decision
  trees with feature subsampling), a feedforward Neural Network (three hidden layers
  with batch normalisation, dropout, and early stopping), XGBoost (regularised
  gradient boosting), and an RBF-kernel SVM with Platt scaling. Cross-validation
  uses Repeated Stratified $k$-Fold with $k=5$ and two repeats (ten folds in total),
  yielding out-of-fold probability vectors for every training sample. These
  out-of-fold probabilities are concatenated into a $460 \times 15$ Level-1 feature
  matrix (five models $\times$ three class probabilities) on which a multinomial
  logistic regression meta-learner is trained via stacked generalisation
  \citep{wolpert1992stacked}. The stacked and simple-averaged ensembles are compared
  on out-of-fold AUC; the better-performing variant is retained for downstream stages.
  The selected ensemble probabilities are then calibrated using both sigmoid and
  isotonic regression \citep{platt1999probabilistic}, with the method that achieves
  the lowest out-of-fold log-loss selected automatically.

  \paragraph{Stage 3 — Escalation Engine and Meta-Model.}
  The escalation engine (\texttt{src/escalation/}) evaluates each patient against
  seven signals derived from the calibrated ensemble output and the pre-imputation
  missingness mask: (i) maximum predicted probability below a confidence threshold
  $\tau_c = 0.70$; (ii) decision margin below $\tau_m = 0.20$; (iii) predictive
  entropy above $\tau_H = 1.00$; (iv) disagreement among the CatBoost, Random Forest,
  and Neural Network top-class predictions; (v) missing-feature fraction exceeding
  $\tau_\text{miss} = 0.30$; (vi) cross-modality mismatch between biomarker and
  cognitive feature directions; and (vii) SHAP attribution variance across the base
  models. Cases triggering signal (v) or (vi) are immediately assigned to the
  Clinician-Mandatory tier; all others are routed to AI-Assisted (one or more signals
  active) or AI-Autonomous (no signals active).

  \paragraph{Stage 4 — Clinician Dashboard.}
  The Streamlit dashboard (\texttt{src/dashboard/app.py}) operates in two modes. In
  \emph{New Patient} mode, a clinician enters or uploads patient biomarker values,
  which are processed through the saved preprocessing artefacts and classified in
  real time. In \emph{Retrospective Review} mode, the clinician selects a data source
  (the 116-patient holdout or the 460-patient OOF set) and browses a ranked patient
  queue. The Patient Detail view presents the predicted class, per-class probability
  bar chart, escalation tier and reasons, a SHAP waterfall plot, biomarker values
  alongside clinical reference ranges, and the five nearest neighbours from the
  training set. A structured feedback form captures the clinician's agreement
  decision, corrected diagnosis, confidence rating (1--5), and free-text notes,
  persisted to a SQLite database.

  \paragraph{Stage 5 — Active-Learning Feedback Loop.}
  The feedback module (\texttt{src/hitl/feedback\_loop.py}) reads clinician
  corrections from the SQLite log, ranks them by escalation risk score, and
  accumulates them in a feedback pool. After each batch of corrections, the full
  ensemble (CatBoost, Random Forest, Neural Network) is retrained on the union of
  the original 460 training samples and the feedback pool. Each corrected sample
  is weighted by $w = 0.5 + c/5$, where $c \in \{1,\ldots,5\}$ is the clinician's
  stated confidence, scaling from $w=0.7$ (confidence 1) to $w=1.0$ (confidence 5).
  In the absence of real clinicians, a configurable simulated-clinician module
  \citep{monarch2021hitl} introduces label noise at a specified accuracy level to
  enable controlled experiments.


  % ===========================================================
  \chapter{Results and Evaluation}
  \label{chap:results}
  % ===========================================================

  This chapter presents quantitative findings from the full evaluation pipeline.
  All reported metrics originate from held-out or out-of-fold data; no training-set
  performance is shown.  The data files cited in each subsection are reproducible
  artefacts stored in \texttt{reports/tables/}.

  % -----------------------------------------------------------
  \section{Baseline Classification Performance}
  \label{sec:baseline-performance}
  % -----------------------------------------------------------

  The selected ensemble's performance on out-of-fold (OOF) predictions
  (\texttt{overall\_metrics\_ensembleselected.csv}) establishes the AI-only baseline
  against which all HITL results are compared:

  \begin{table}[htbp]
    \centering
    \caption{Ensemble OOF classification metrics (AI-only baseline).}
    \label{tab:baseline}
    \begin{tabular}{lS[table-format=1.3]}
      \toprule
      {Metric} & {Value} \\
      \midrule
      Accuracy            & 0.730 \\
      Macro Precision     & 0.733 \\
      Macro Recall        & 0.739 \\
      Macro F1            & 0.728 \\
      Weighted F1         & 0.733 \\
      Macro AUC (OvR)     & 0.902 \\
      Log-loss            & 0.553 \\
      ECE                 & 0.044 \\
      \bottomrule
    \end{tabular}
  \end{table}

  Per-class analysis (\texttt{per\_class\_metrics\_ensembleselected.csv}) reveals
  strongest discrimination for AD (AUC\,=\,0.955) and SCD (AUC\,=\,0.938), with
  MCI as the hardest class (AUC\,=\,0.818; F1\,=\,0.605).  This asymmetry is
  clinically expected: the intermediate stage sits between cognitively normal and
  demented, producing the most ambiguous biomarker profiles.

  Bootstrap confidence intervals (\texttt{bootstrap\_ci\_ensembleselected.csv})
  confirm that the reported point estimates are stable under resampling.  The macro
  AUC 95\% CI is [0.879, 0.924], indicating that even in the worst case the model
  substantially exceeds chance across all three classes.

  % -----------------------------------------------------------
  \section{Meta-Model Selection and Validation}
  \label{sec:meta-model}
  % -----------------------------------------------------------

  The meta-model that converts the seven escalation signals into a scalar risk
  score was selected by cross-validated AUC over three candidate architectures
  (\texttt{step2\_meta\_model\_comparison.csv}):

  \begin{table}[htbp]
    \centering
    \caption{Meta-model candidate comparison (5-fold CV on OOF data).}
    \label{tab:meta-model}
    \begin{tabular}{lSS}
      \toprule
      {Model} & {CV AUC} & {CV Accuracy} \\
      \midrule
      Logistic Regression   & 0.759 & 0.665 \\
      Random Forest         & 0.758 & 0.720 \\
      Gradient Boosting     & 0.728 & 0.726 \\
      \bottomrule
    \end{tabular}
  \end{table}

  Logistic Regression was selected because it achieved the highest AUC (0.759),
  and its monotonic, interpretable decision boundary is preferable in a clinical
  context where clinicians must understand \emph{why} a case was escalated.
  Bootstrap validation of the selected model
  (\texttt{step2\_meta\_bootstrap\_metrics.csv}, $B = 1000$) yielded a mean AUC of
  $0.758 \pm 0.023$ with 95\% CI $[0.713, 0.801]$, confirming that the meta-model's
  discriminatory ability is not an artefact of a single split.

  % -----------------------------------------------------------
  \section{Escalation Signal Ablation}
  \label{sec:ablation}
  % -----------------------------------------------------------

  To assess the marginal contribution of each escalation signal, an ablation study
  was conducted in which each of the seven risk features was removed in turn and the
  meta-model was retrained (\texttt{ablation\_study.csv}):

  \begin{table}[htbp]
    \centering
    \caption{Ablation study: effect of removing each escalation signal on
    meta-model AUC.}
    \label{tab:ablation}
    \begin{tabular}{lSS}
      \toprule
      {Removed feature} & {$\Delta$ AUC} & {$\Delta$ Accuracy} \\
      \midrule
      None (full model)              & {---}   & {---}  \\
      \texttt{risk\_1\_inv\_conf}    & -0.001  & 0.000  \\
      \texttt{risk\_2\_margin}       & -0.018  & -0.007 \\
      \texttt{risk\_3\_entropy}      & +0.002  & +0.020 \\
      \texttt{risk\_4\_disagreement} &  0.000  & +0.002 \\
      \texttt{risk\_5\_missing}      & -0.000  &  0.000 \\
      \texttt{risk\_6\_critical}     & -0.000  &  0.000 \\
      \texttt{risk\_7\_multimodal}   & +0.003  & +0.009 \\
      \bottomrule
    \end{tabular}
  \end{table}

  The decision-margin signal (\texttt{risk\_2\_margin}) carries the largest
  individual predictive weight: its removal produces the steepest AUC decline
  ($-0.018$).  By contrast, inverse confidence (\texttt{risk\_1\_inv\_conf}) and
  the two missingness signals contribute negligibly in isolation.  Removing entropy
  or multimodal mismatch \emph{increases} the metric slightly, suggesting partial
  redundancy with margin; the full feature set is retained because the combined
  signal set improves threshold sensitivity robustness (Section~\ref{sec:threshold}).

  % -----------------------------------------------------------
  \section{Threshold Selection and Sensitivity}
  \label{sec:threshold}
  % -----------------------------------------------------------

  The escalation threshold $\tau^*$ governs the trade-off between review rate
  (clinician workload) and policy accuracy.  Following the procedure described in
  Section~\ref{sec:architecture}, $\tau^*$ was selected to maximise expected policy
  accuracy under a uniform distribution of assumed clinician accuracy
  $h \sim \mathcal{U}[0.85, 0.95]$, subject to a review-budget cap of 40\%.

  The selected threshold is $\tau^* = 0.55$, which routes 39.8\% of cases for human
  review.  Table~\ref{tab:threshold-sensitivity} shows that the threshold choice is
  \emph{stable}: the same $\tau^*$ is optimal for every tested value of $h$.

  \begin{table}[htbp]
    \centering
    \caption{Threshold sensitivity to assumed clinician accuracy ($\tau^* = 0.55$
    throughout).}
    \label{tab:threshold-sensitivity}
    \begin{tabular}{SSSS}
      \toprule
      {$h$} & {Policy accuracy} & {AI-only accuracy} & {Accuracy gain} \\
      \midrule
      0.85 & 0.860 & 0.730 & 0.129 \\
      0.88 & 0.872 & 0.730 & 0.141 \\
      0.90 & 0.880 & 0.730 & 0.149 \\
      0.92 & 0.888 & 0.730 & 0.157 \\
      0.95 & 0.900 & 0.730 & 0.169 \\
      \bottomrule
    \end{tabular}
  \end{table}

  At the assumed midpoint $h = 0.90$, the HITL policy achieves 88.0\% accuracy ---
  a $+14.9$ percentage-point gain over the AI-only baseline.  Even in the worst case
  ($h = 0.85$), the gain is $+12.9$ pp, and the policy accuracy (86.0\%) remains
  substantially above the AI-only level.  The best-case policy accuracy is 90.0\%.

  % -----------------------------------------------------------
  \section{AI-Only versus HITL Performance}
  \label{sec:hitl-comparison}
  % -----------------------------------------------------------

  The core experimental comparison evaluates the complete HITL pipeline against
  the AI-only baseline on out-of-fold predictions ($n = 460$), using the selected
  threshold $\tau^* = 0.55$ and simulated clinician accuracy $h = 0.90$
  (\texttt{statistical\_tests.csv}).

  \begin{table}[htbp]
    \centering
    \caption{AI-only vs HITL performance with statistical inference.}
    \label{tab:hitl-main}
    \begin{tabular}{lcc}
      \toprule
      & AI-only & HITL ($\tau^* = 0.55$, $h = 0.90$) \\
      \midrule
      Accuracy      & 0.730 [0.689, 0.770] & 0.889 [0.861, 0.917] \\
      Macro F1      & 0.728 [0.687, 0.766] & 0.888 [0.859, 0.917] \\
      \midrule
      \multicolumn{3}{l}{\textit{Paired comparison}} \\
      Accuracy gain &  \multicolumn{2}{c}{+0.159 [0.124, 0.196]} \\
      Macro F1 gain &  \multicolumn{2}{c}{+0.161 [0.126, 0.197]} \\
      McNemar $\chi^2$ & \multicolumn{2}{c}{$60.99$ ($p = 5.8 \times 10^{-15}$)} \\
      Matched OR    & \multicolumn{2}{c}{$13.17$} \\
      \bottomrule
    \end{tabular}
  \end{table}

  The HITL pipeline improves accuracy by $+15.9$ percentage points (95\% CI
  $[12.4, 19.6]$) and macro F1 by $+16.1$ pp (95\% CI $[12.6, 19.7]$).
  McNemar's test confirms the difference is statistically significant at
  $p < 10^{-14}$, with a matched odds ratio of 13.17 ($b = 6$ corrections
  worsened by HITL, $c = 79$ corrections improved), meaning HITL is approximately
  13 times more likely to correct an AI error than to introduce one.

  \subsection{Per-Class Breakdown}

  Table~\ref{tab:perclass} disaggregates the improvement by diagnostic class
  (\texttt{hitl\_perclass\_comparison.csv}):

  \begin{table}[htbp]
    \centering
    \caption{HITL accuracy gain by diagnostic class.}
    \label{tab:perclass}
    \begin{tabular}{lSSSS}
      \toprule
      {Class} & {$n$} & {AI-only acc.} & {HITL acc.} & {Gain (pp)} \\
      \midrule
      SCD & 154 & 0.818 & 0.931 & +11.3 \\
      MCI & 156 & 0.564 & 0.767 & +20.3 \\
      AD  & 150 & 0.813 & 0.945 & +13.1 \\
      \bottomrule
    \end{tabular}
  \end{table}

  The largest improvement occurs for MCI ($+20.3$ pp), the class that the AI-only
  model found most difficult.  This is the intended behaviour of the escalation
  engine: cases near the SCD--MCI and MCI--AD boundaries produce high-entropy,
  low-margin predictions, which trigger escalation to a human reviewer who resolves
  the ambiguity.  SCD and AD, although already well-classified, still gain $+11$--$13$
  pp from targeted human review.

  % -----------------------------------------------------------
  \section{Escalation Quality Analysis}
  \label{sec:escalation-quality}
  % -----------------------------------------------------------

  The quality of the escalation decision itself --- does the meta-model
  \emph{correctly} identify cases that the AI got wrong? --- is assessed via
  precision, recall, and F1 at the operating point $\tau^* = 0.55$
  (\texttt{escalation\_quality\_operating\_point.csv}):

  \begin{table}[htbp]
    \centering
    \caption{Escalation quality at the operating point ($\tau^* = 0.55$).}
    \label{tab:escalation-quality}
    \begin{tabular}{lS}
      \toprule
      {Metric} & {Value} \\
      \midrule
      Review rate           & 0.398 \\
      Precision (error detection) & 0.475 \\
      Recall (error detection)    & 0.702 \\
      F1                          & 0.567 \\
      Specificity                 & 0.714 \\
      Missed-error rate           & 0.298 \\
      Unnecessary escalation rate & 0.525 \\
      \bottomrule
    \end{tabular}
  \end{table}

  The meta-model captures 70.2\% of the AI's errors at a review rate of 39.8\%.
  Precision is moderate (47.5\%), meaning roughly half of escalated cases were
  already correctly classified; these ``unnecessary'' escalations impose clinician
  cost but do not harm accuracy.  The budget-sweep analysis
  (\texttt{escalation\_quality\_budget\_summary.csv}) shows how this trade-off
  evolves: at a 10\% budget the recall drops to 16.9\%, while at 50\% it rises to
  71.8\% with diminishing returns beyond that point.

  % -----------------------------------------------------------
  \section{Cost--Benefit Analysis}
  \label{sec:cost}
  % -----------------------------------------------------------

  To translate raw accuracy gains into clinical utility, a cost model assigns
  asymmetric misclassification costs: missed AD diagnoses carry higher penalties
  than over-referral of SCD cases, reflecting the clinical reality that a delayed
  dementia diagnosis has graver consequences than an unnecessary specialist
  appointment.  Results are drawn from \texttt{clinical\_cost\_summary.csv}:

  \begin{table}[htbp]
    \centering
    \caption{Clinical cost comparison at $\tau^* = 0.54$ (selected by cost
    objective).}
    \label{tab:cost}
    \begin{tabular}{lSSS}
      \toprule
      {$h$} & {AI cost} & {HITL cost} & {Cost reduction} \\
      \midrule
      0.85 & 209.0 & 89.5  & 119.5 \\
      0.90 & 209.0 & 77.0  & 132.0 \\
      0.95 & 209.0 & 70.5  & 138.5 \\
      \bottomrule
    \end{tabular}
  \end{table}

  The HITL pipeline reduces the clinically weighted cost by 57--66\% depending on
  assumed clinician accuracy.  Even with imperfect human reviewers ($h = 0.85$),
  cost-weighted accuracy improves from 90.9\% to 96.1\%, supporting the claim that
  the framework improves not only statistical performance but also clinical utility.

  % -----------------------------------------------------------
  \section{Conformal Prediction and Uncertainty Quantification}
  \label{sec:conformal}
  % -----------------------------------------------------------

  Conformal prediction provides distribution-free coverage guarantees: for a
  user-specified miscoverage level $\alpha$, the procedure outputs a prediction
  \emph{set} that contains the true class with probability at least $1 - \alpha$
  (\texttt{conformal\_prediction.csv}).

  \begin{table}[htbp]
    \centering
    \caption{Conformal prediction coverage and set-size statistics.}
    \label{tab:conformal}
    \begin{tabular}{SSSSS}
      \toprule
      {$\alpha$} & {Target cov.} & {Empirical cov.} & {Avg.\ set size}
        & {Singleton frac.} \\
      \midrule
      0.05 & 0.95 & 1.00 & 2.61 & 0.000 \\
      0.10 & 0.90 & 1.00 & 2.54 & 0.000 \\
      0.15 & 0.85 & 1.00 & 2.42 & 0.017 \\
      0.20 & 0.80 & 1.00 & 2.34 & 0.026 \\
      \bottomrule
    \end{tabular}
  \end{table}

  Empirical coverage is 100\% at every tested $\alpha$, exceeding the nominal
  targets.  This conservatism --- desirable in a clinical setting where
  under-coverage would miss true diagnoses --- comes at the cost of large prediction
  sets (mean 2.3--2.6 out of 3 classes).  Singleton sets, which would be maximally
  informative to a clinician, appear in fewer than 3\% of cases.

  The over-coverage indicates that the calibrated ensemble's softmax probabilities
  are already well-separated for the majority of cases; the conformal threshold
  $\hat{q}$ is pushed high enough to guarantee marginal coverage, which means
  nearly all classes pass the threshold for most patients.  Future work could
  explore conditional (class-wise or group-wise) conformal methods to tighten the
  sets while maintaining coverage.

  \subsection{Uncertainty Decomposition}

  Predictive uncertainty was decomposed into aleatoric and epistemic components
  using the ensemble's disagreement structure
  (\texttt{uncertainty\_decomposition.csv}):

  \begin{table}[htbp]
    \centering
    \caption{Uncertainty decomposition across test samples.}
    \label{tab:uncertainty}
    \begin{tabular}{lSSS}
      \toprule
      {Component} & {Mean} & {Std} & {Median} \\
      \midrule
      Total uncertainty     & 0.674 & 0.222 & 0.736 \\
      Aleatoric (data)      & 0.644 & 0.216 & 0.690 \\
      Epistemic (model)     & 0.030 & 0.023 & 0.023 \\
      Epistemic fraction    & 0.047 & 0.033 & 0.038 \\
      \bottomrule
    \end{tabular}
  \end{table}

  Total uncertainty is dominated by aleatoric uncertainty (mean epistemic
  fraction\,=\,4.7\%), implying that most predictive difficulty arises from
  inherent overlap between the three diagnostic classes rather than from
  insufficient training data or model capacity.  This finding aligns with clinical
  knowledge: the SCD--MCI boundary is biologically gradual, not a sharp cut-off.

  The low epistemic fraction also suggests that the five-model ensemble is well
  diversified; adding further base learners would yield diminishing returns, whereas
  collecting additional biomarker modalities (e.g.\ neuroimaging, genetic markers)
  could reduce the dominant aleatoric component.

  % -----------------------------------------------------------
  \section{Fairness Analysis}
  \label{sec:fairness}
  % -----------------------------------------------------------

  Algorithmic fairness was evaluated across two protected attributes --- gender
  and age group --- using demographic parity difference and equalized odds
  difference (\texttt{fairness\_summary.csv}):

  \begin{table}[htbp]
    \centering
    \caption{Fairness metrics by protected attribute.}
    \label{tab:fairness}
    \begin{tabular}{lSS}
      \toprule
      {Attribute} & {Demographic parity diff.} & {Equalized odds diff.} \\
      \midrule
      Gender    & 0.009 & 0.087 \\
      Age group & 0.163 & 0.147 \\
      \bottomrule
    \end{tabular}
  \end{table}

  Gender-based disparities are small: the escalation rate differs by less than one
  percentage point between groups, and equalized odds difference remains below 0.10.
  Age-group disparities are materially larger (demographic parity
  difference\,=\,0.163), indicating that the escalation engine routes patients of
  different age bins at substantially different rates.  This is partly expected ---
  AD prevalence increases with age, producing genuine base-rate differences --- but
  it warrants monitoring.  If deployed, age-stratified threshold calibration or
  post-hoc equalized-odds adjustments could mitigate the disparity without
  significantly reducing accuracy.

  % -----------------------------------------------------------
  \section{Active-Learning Feedback Loop}
  \label{sec:active-learning}
  % -----------------------------------------------------------

  The feedback loop was evaluated using simulated clinicians at three accuracy
  levels ($h \in \{0.85, 0.90, 0.95\}$).  After each round, the top-10
  highest-risk-score cases were ``reviewed'' by the simulated clinician and their
  (possibly noisy) labels were added to the training set
  (\texttt{active\_learning\_summary.csv}):

  \begin{table}[htbp]
    \centering
    \caption{Active-learning feedback: accuracy after one round of 10 corrections.}
    \label{tab:active-learning}
    \begin{tabular}{lSSSS}
      \toprule
      {$h$} & {Round 0 acc.} & {Round 1 acc.} & {$\Delta$ acc.\ (pp)}
        & {$\Delta$ F1 (pp)} \\
      \midrule
      0.85 & 0.733 & 0.741 & +0.86 & +1.50 \\
      0.90 & 0.733 & 0.741 & +0.86 & +0.97 \\
      0.95 & 0.733 & 0.741 & +0.86 & +0.97 \\
      \bottomrule
    \end{tabular}
  \end{table}

  Even a single round of 10 clinician-corrected cases produces a measurable
  improvement ($+0.86$ pp accuracy) regardless of clinician accuracy, because the
  escalation engine correctly prioritises the cases the model is most uncertain
  about.  The accuracy gain is identical across clinician levels at round 1 because
  10 corrections are too few for the noise difference ($h = 0.85$ vs $0.95$) to
  manifest; longer feedback sequences would be expected to separate the curves.

  The modest magnitude of the gain reflects the small feedback pool size relative to
  the 460-sample training set.  In a real deployment where hundreds of cases
  accumulate over months, the compounding effect of targeted corrections on the
  hardest boundary cases (particularly MCI) is expected to be substantially larger.

  % -----------------------------------------------------------
  \section{Summary of Evidence}
  \label{sec:results-summary}
  % -----------------------------------------------------------

  Table~\ref{tab:evidence-summary} consolidates the key findings across all
  evaluation dimensions:

  \begin{table}[htbp]
    \centering
    \caption{Summary of experimental evidence.}
    \label{tab:evidence-summary}
    \begin{tabular}{p{3.8cm}p{8.5cm}}
      \toprule
      \textbf{Dimension} & \textbf{Key finding} \\
      \midrule
      HITL accuracy gain & $+15.9$ pp over AI-only ($p < 10^{-14}$);
        largest for MCI ($+20.3$ pp). \\
      Threshold robustness & Same $\tau^*$ is optimal across
        $h \in [0.85, 0.95]$; review rate 39.8\%. \\
      Escalation quality & 70.2\% of AI errors captured at
        39.8\% review rate. \\
      Clinical cost & HITL reduces cost by 57--66\%
        vs AI-only. \\
      Conformal coverage & 100\% empirical coverage at all $\alpha$;
        conservative but safe. \\
      Uncertainty structure & 95\% aleatoric, indicating class overlap
        is the dominant challenge. \\
      Fairness & Gender disparity minimal; age-group disparity
        requires monitoring. \\
      Active learning & $+0.86$ pp from a single round of 10 corrections. \\
      \bottomrule
    \end{tabular}
  \end{table}

  Taken together, these results demonstrate that the HITL framework consistently
  improves classification performance, clinical utility, and uncertainty
  characterisation relative to the AI-only baseline.  The principal remaining risks
  are (i) age-group fairness disparity in escalation rates, and (ii) the need
  for real clinician validation to replace the simulated-clinician assumption.


  % ===========================================================
  %  BIBLIOGRAPHY
  % ===========================================================
  \bibliographystyle{plainnat}
  \bibliography{refs}

  \end{document}
